\documentclass[11pt]{article}
\usepackage[T1]{fontenc}
\usepackage{lmodern,amsmath,amssymb,amsthm}
\usepackage[margin=1in]{geometry}
\usepackage{microtype}
\usepackage[colorlinks=true,urlcolor=blue,linkcolor=blue]{hyperref}
\newtheorem{corollary}{Corollary}
\newtheorem{lemma}[corollary]{Lemma}
\newcommand{\E}{\mathbb E}
\newcommand{\Prb}{\mathbb P}
\title{A lower bound of 3.671741 for the random 3-SAT threshold}
\author{Fedor Vorobyev}
\date{Technical note, version 1.0\\7 October 2026}
\hypersetup{pdftitle={A lower bound of 3.671741 for the random 3-SAT threshold},pdfauthor={Fedor Vorobyev}}
\begin{document}
\maketitle
\begin{abstract}
Combining the spatial-coupling lower bound on the zero-violation-density
threshold with the new linear violation robustness theorem yields
\[
 \alpha_3\ge
 A:=\inf_{0\le x<1}\frac{4(1-x+x\log x)}{(1-x)^3}
 >3.671741.
\]
The first input is due to Achlioptas, Hassani, Macris, and Urbanke;
the second is Proposition 4.6 of OpenAI's \emph{Computing the Random
3-SAT Threshold}. We give the transfer argument and a finite exact
rational certificate for the numerical inequality. The constant $A$
comes from the earlier spatial-coupling analysis; the deduction here
transfers that bound to exact satisfiability.
\end{abstract}

\section{Statement and the missing implication}
Let $\Phi(n,m)$ have $n$ Boolean variables and $m$ independent clauses,
each using three uniformly chosen distinct variables with independent
fair signs. Whole clauses may repeat. Let $\alpha_3$ be its limiting
satisfiability threshold, whose existence is established in~\cite{Carenini}
and also in~\cite[Theorem 4.4]{OpenAI}, and let
\[
 V(F)=\min_{\sigma}\#\{\text{clauses of }F\text{ violated by }\sigma\}.
\]
\begin{corollary}\label{thm:main}
Taking the spatial-coupling theorem of \cite{AHMU,HMU} and the robustness
theorem of \cite{OpenAI} as inputs,
\[
 \alpha_3\ge A>\frac{3671741}{10^6},
 \qquad
 A=\inf_{0\le x<1}\frac{4(1-x+x\log x)}{(1-x)^3},
\]
where $0\log0=0$. In particular, for every fixed
$0\le a\le3671741/10^6$, $\Phi(n,\lfloor an\rfloor)$ is satisfiable with
probability tending to one.
\end{corollary}

Together with the upper-bound theorem of~\cite{Vorobyev}, this gives
\[
 \boxed{3.671741<\alpha_3\le4.268.}
\]

The spatial-coupling result concerns vanishing minimum violation
\emph{density}, rather than exact satisfaction. The new robustness
proposition supplies the necessary connection. Threshold existence or
computability alone is not substituted for that proposition.

\section{Two probabilistic inputs}
\paragraph{Spatial coupling.}
Achlioptas, Hassani, Macris, and Urbanke establish a lower bound on the
rough threshold: at every fixed density $a<A$, the limiting expected
minimum number of violations per variable is zero~\cite{AHMU}.
Their full manuscript, Theorem 4 and equations (90) and (92), characterizes the
coupled unit-clause threshold using
\[
 \Phi_{\rm ind}(\ell,a,3)
 =2-\ell\bigl(1-\log(\ell/2)\bigr)
       -\frac a2(1-\ell/2)^3.
\]
Setting $\ell=2x$ gives the expression for $A$ above. The passage to
ordinary formulas uses the asymptotic energy interpolation of
Hassani, Macris, and Urbanke~\cite[Theorems 1--2]{HMU}.
One takes the number of variables per position to infinity, then
the chain length to infinity at fixed coupling width, and finally
the width to infinity. We use this limiting statement, not a
finite-chain equality of thresholds.

\paragraph{Linear violation robustness.}
In the relaxed Poisson model, variable positions inside each clause
are independently sampled, and the number of clauses is Poisson with
mean $an$. Proposition 4.6 of the OpenAI preprint~\cite{OpenAI} states that, for every fixed
$a>\alpha_3$, there exists $\varepsilon(a)>0$ such that
\begin{equation}\label{eq:robust}
 \Prb\{V(F^{\rm rep}_{n,\operatorname{Poi}(an)})
                   >\lfloor\varepsilon(a)n\rfloor\}\longrightarrow1.
\end{equation}
Theorem 4.4 of that paper identifies the threshold in this model
with the proper-clause threshold. The robustness proof uses its
free-deletion estimates; it does not assume the spatial-coupling
numerical bound or the predicted threshold value.

\section{Transferring zero violation density}
Adding, deleting, or replacing one clause changes $V$ by at most one.
This gives the needed model transfers without an assertion about
finite-size satisfiability probabilities.

Write $F^{\rm rep}_{n,m}$ for fixed-count clauses with independently
sampled variable positions, and $F^{\rm prop}_{n,m}$ for proper clauses.
Resample every nonproper clause independently from the proper law.
Then, with $m=\lfloor an\rfloor$,
\begin{equation}\label{eq:proper}
 |\E V(F^{\rm rep}_{n,m})-\E V(F^{\rm prop}_{n,m})|
 \le m\left[1-(1-1/n)(1-2/n)\right]\le3m/n.
\end{equation}
For an independent $M\sim\operatorname{Poi}(an)$, use one common
infinite relaxed clause stream. Cauchy--Schwarz gives
\begin{equation}\label{eq:poisson}
 |\E V(F^{\rm rep}_{n,M})-\E V(F^{\rm rep}_{n,m})|
 \le\E|M-m|\le\sqrt{an+1}.
\end{equation}
After division by $n$, both errors vanish. Thus the zero-energy
conclusion from spatial coupling transfers to the model of
\eqref{eq:robust}.

Fix $a<A$. The spatial-coupling conclusion is
\[
 \frac1n\E V(F^{\rm rep}_{n,\operatorname{Poi}(an)})\longrightarrow0.
\]
If $a>\alpha_3$, equation~\eqref{eq:robust} instead gives
\[
 \liminf_{n\to\infty}\frac1n
       \E V(F^{\rm rep}_{n,\operatorname{Poi}(an)})
 \ge\varepsilon(a)>0,
\]
a contradiction. Therefore $a\le\alpha_3$ for every $a<A$, proving
$\alpha_3\ge A$. Only a numerical lower bound for $A$ remains.

\section{An exact rational certificate}
Set
\[
 a_0=\frac{3671741}{10^6},\qquad N=128,\qquad
 y=\frac{883413964}{10^9},\qquad \eta=\frac9{10^8}.
\]
For $t=1-x\in[0,1]$, the entropy series is
\[
 1-x+x\log x
 =\sum_{k=2}^{\infty}\frac{t^k}{k(k-1)}.
\]
At $t=1$ the identity follows by continuity or telescoping.
Define the rational polynomial
\[
 P(t)=\sum_{k=2}^{128}\frac{4t^{k-2}}{k(k-1)}-a_0t.
\]
Every coefficient of $P''$ is nonnegative, so $P$ is convex on $[0,1]$.
Its tangent at $y$ therefore gives, for every $t\in[0,1]$,
\[
 P(t)\ge P(y)+P'(y)(t-y)\ge P(y)-|P'(y)|.
\]
Exact rational arithmetic proves
\begin{equation}\label{eq:certificate}
 \boxed{P(y)-|P'(y)|>\frac9{10^8}.}
\end{equation}
The accompanying standard-library Python program
\texttt{certify\_3671741.py} evaluates $P$ and $P'$ both by Horner's
rule and by direct summation, and checks \eqref{eq:certificate} using
integer comparisons. For orientation only, the rational values are
approximately
\[
 P(y)=9.3928406844\cdot10^{-8},\qquad
 P'(y)=3.7588568378\cdot10^{-9}.
\]
These decimal values are not used by the verifier.

The omitted series terms are nonnegative. Thus for $0<t\le1$,
\begin{align*}
 \frac{4(1-x+x\log x)}{(1-x)^3}-a_0
 &=\frac1t\left(
       \sum_{k=2}^{\infty}\frac{4t^{k-2}}{k(k-1)}-a_0t\right)\\
 &\ge\frac{P(t)}t>\eta.
\end{align*}
Taking the infimum proves $A\ge a_0+\eta>a_0$, completing
Corollary~\ref{thm:main}. In particular the theorem's strict inequality
does not depend on interpreting a rounded decimal as an exact bound.

The stationary equation for the interior minimizer is
\[
 (1+2x)\log x+3(1-x)=0.
\]
Its numerical solution gives $x\approx0.11658603275812085$ and
$A\approx3.671741106632344$. These approximations describe the exact
variational constant; the rational certificate above is sufficient
for the stated lower bound.

\section{Attribution and verification}
The constant $A$ is inherited from the spatial-coupling work of
Achlioptas, Hassani, Macris, and Urbanke. The contribution of this note
is the deduction of an exact-satisfiability lower bound using the
violation-robustness proposition of~\cite{OpenAI}. The upper bound
$4.268$ is a separate result~\cite{Vorobyev} and is not an input to
the lower-bound argument.

The probabilistic results are used as cited theorems. The accompanying
verifier checks the scalar inequality using exact rational arithmetic;
it does not verify the proofs of those theorems. This note has no
complete Lean formalization.

\paragraph{Use of language models.}
OpenAI language models assisted with identifying the combination of
results, checking the deduction, developing the arithmetic certificate,
and drafting this note. Additional checks were performed in separate
model runs; these do not constitute independent human peer review.
The author directed the work and is responsible for the final text.

\begin{thebibliography}{9}
\bibitem{AHMU}
D. Achlioptas, S. H. Hassani, N. Macris, and R. Urbanke.
\emph{Bounds for Random Constraint Satisfaction Problems via Spatial
Coupling}. SODA 2016, pp. 469--479.
\href{https://doi.org/10.1137/1.9781611974331.ch35}{doi:10.1137/1.9781611974331.ch35}.
The authors' full manuscript is available at
\url{https://di.uoa.gr/~optas/papers/spatial.pdf}; Theorem 4 and
equations (90) and (92) are in the appended full version, pp. 29--30.
\bibitem{HMU}
S. H. Hassani, N. Macris, and R. Urbanke.
\emph{Threshold Saturation in Spatially Coupled Constraint Satisfaction
Problems}. Journal of Statistical Physics 150 (2013), pp. 807--850.
\url{https://arxiv.org/abs/1112.6320v2}.
\bibitem{OpenAI}
OpenAI. \emph{Computing the Random 3-SAT Threshold}, preprint dated
27 September 2026. Theorem 4.4, p. 11, and Proposition 4.6, pp. 13--14.
\href{https://github.com/openai/math/tree/adc7f1241b42e322a6451854ab7e4b4c146bf78a/preprints/Computing-the-Random-3-SAT-Threshold-September-27-2026}{Manuscript and source at the cited revision}.
\bibitem{Carenini}
G. Carenini. \emph{A polynomial scaling window for random $k$-SAT
and a proof of the satisfiability conjecture}. ECCC TR26-229,
5 October 2026. \url{https://eccc.weizmann.ac.il/report/2026/229/}.
\bibitem{Vorobyev}
F. Vorobyev. \emph{An upper bound of 4.268 for the random 3-SAT
satisfiability threshold}, October 2026.
\url{https://github.com/yamelton/sat-upper-4268}.
\end{thebibliography}
\end{document}
